\paragraph{Conjugate-coordinate class: fibre admissibility and the Lorentzian \(i\epsilon\) window (conditional; bounded no-go).}
Keep the setting and notation of the preceding paragraph (Poisson limit, leaf \(\mathbb{CP}^{1,2}\), \(T^a=M^{a4}\), \(\Box=\eta_{ab}D^aD^b\), \(H_V\) with \(\mu^2=-3\), \(C_{42}\), \(C_{41}\), \(C_{31}\), \(\sigma(\cdot)\), \(\Delta_\kappa\)); write \(S^2=2C_{41}-C_{42}\), whose symbol \(\beta q_f\) is purely vertical.

\emph{Fibre cutoffs.} In the Poisson limit the fibre of the leaf is a compact two-sphere and its harmonic number, equivalently the centre variable \(\zeta\), is unbounded; this is the sense of ``unbounded fibre'' in the preceding paragraphs. Two bounds on the harmonic number are in the literature. (a) \(S^2=2s(s+1)\) on the sectors \(\mathcal C^s\) with \(s\le n\), from the \((n+1)\)-dimensional representation of the stabiliser \(SU(2)_L\) on the lowest-weight state~\cite{Steinacker2019NoGhost}; the classical counterpart is the radius of the fibre as a coadjoint orbit of \(SU(2)_L\), the self-dual part of \(M^{ab}\) on the \(\eta\)-orthogonal complement of \(y\): at \(10\) orbit points exactly one of the two radii vanishes and the other equals \(\ell\), independent of \(y\). (b) A local bound \(s\le m=r^{-1}x^0\) from the uncertainty of the coherent states, which the same group states to correct earlier statements on the higher-spin cutoff~\cite{MantaSteinacker2025}; it grows with the point and is not the area bound of (a), and the two are not reconciled here. Neither is a sharp cutoff of commuting matrix-built operators.

\emph{Admissible cone.} For admissible modes, the principal series of \(SO(4,1)\), one has \(r^{-2}\Box_H>s^2+s+9/4\) with \(\Box_H=-C_{42}+C_{41}\) in the convention of~\cite{Steinacker2019NoGhost}, equivalent to \(-C_{42}>9/2\) in that convention (stated there as plausible, not proved, for the \(SO(4,2)\) content of \(\operatorname{End}(\mathcal H_n)\)). With \(\Box_H\) horizontal, \(S^2\) vertical and \(\sigma(C_{42})=-\alpha q_b+\beta q_f\) of the preceding paragraph, the semiclassical content of this bound is the cone \(-\sigma(C_{42})\ge0\), that is \(\alpha q_b\ge\beta q_f\); the Casimirs here and there differ by a sign at the symbol level. Conditional on this identification, a cutoff of \(C_{42}\) restricted to the cone together with a bound \(\beta q_f\le Q\) is a compact region, so a function of the central operator \(C_{42}\) regularises once the fibre is bounded. In that case \(\Delta_\kappa\) is not needed: for fixed \(Q\) and \(\LUV^2\to\infty\) the window \(\{-\sigma(\Delta_\kappa)\le\LUV^2\}\) has its vertical term bounded by \(\beta(1-\kappa)Q\), so \(\kappa\) only rescales \(\LUV^2\) by \(2\kappa-1\) at leading order and the divergence of the window volume at the ends of the interval is absent. The cutoff function and the base scale remain external: the quantisation map is \(SO(4,2)\)-equivariant and so acts on each irreducible component by a scalar, but this multiplier is not computed, and neither \(\chi\) nor the base scale is fixed. The window uses \(M^{AB}\), which are background data and not dynamical matrices, so the regulated trace is not shown to be a diffeomorphism-covariant functional of \(X^a\).

\emph{Lorentzian trace in a window.} Define the one-loop sum as \(\sum\ln(\lambda-i\epsilon)=-\int_0^\infty\frac{\mathrm ds}{s}\operatorname{Tr}\bigl[\chi\,e^{-is(H_V-i\epsilon)}\bigr]\) up to a scale-dependent constant, with \(\chi\) a function of the commuting operator \(\Delta_\kappa\); the contour is the proper time of \(H_V\) with the factor \(e^{-\epsilon s}\), and uses no metric. In a window of compact support (or a Gaussian) the number of modes per unit volume is finite, so the trace is finite for every \(s\), and \(\sum\ln(\lambda-i\epsilon)\) is a finite product over the window whether or not the Schwinger integral converges. For \(\chi=e^{-\sigma_W/\LUV^2}\), \(\sigma_W=-\sigma(\Delta_\kappa)\), the leading-symbol density is the Gaussian \((4\pi)^{-3}\LUV^6\det(\sigma_W)^{-1/2}\prod_{k=1}^{4}(1+is\LUV^2\lambda_k)^{-1/2}\), with \(\lambda_k\) the nonzero generalised eigenvalues of \((\sigma(\Box),\sigma_W)\); the kernel of \(\sigma(\Box)\) is not the vertical plane (checked at eight points). The mass-type moment is \(\langle(\sigma-i0)^{-1}\rangle=i\int_0^\infty\mathrm ds\prod_k(1+is\lambda_k)^{-1/2}\), whose real part is the principal value. For \(\lambda=(-c,a,a,a)\) the real part is \(2/(a+c)\) (closed form, checked), and for \(4000\) samples with \(c\in[0.1,10]\), \(a_k\in[10^{-3},10^{3}]\) it is positive. At eight orbit points (\(\kappa=\tfrac34\), \((y^4)^2\) from \(0.08\) to \(144\)) it is \(1.1\) to \(2.3\) times the value for the Euclidean contraction \(\eta\to\delta\), with the same sign, and the imaginary part is positive.

\emph{Flux.} At the same points the Lorentzian centre flux \(F_{ab}(\zeta)\) has \(\eta F\) with one real pair \(\pm a\) and one imaginary pair \(\pm ib\), \(a,b>0\) at all sampled \(\zeta\), and \(F_{ab}F^{ab}\) takes both signs (negative for \((y^4)^2\lesssim0.4\), positive above). The vector block therefore has Zeeman eigenvalues \(\pm2b\) (phases in \(e^{-is(H_V-i\epsilon)}\)) and \(\pm2ia\) (a factor \(e^{\pm2sa}\)). The Schwinger integrand of the vector sector then grows for the electric pair, bounded by \(e^{2sa}\) with \(a\) linear in the cutoff, while \(\ln\det(\sigma+Z)=\ln(\sigma^2-4b^2)+\ln(\sigma^2+4a^2)\) is finite. The exponential growth of the Euclidean contraction is absent order by order in the flux (each order is a polynomial in \(s\) times \(e^{-is\sigma}\)) and for the magnetic pair, and is present for the electric pair as a property of the integral representation and not of the determinant. The order-\(F^2\) term is proportional to \(F_{ab}F^{ab}\), so its sign is not definite on the leaf.

\emph{Invariant sector, exactly.} On fibre-constant functions \(\Box=-(\Delta_{H^4}-\Delta_{H^3})\) with \(H^4=\mathbb R_\tau\times H^3\), metric \(\mathrm d\tau^2+\cosh^2\tau\,\mathrm dH_3^2\), \(y^4=\ell\sinh\tau\), and the diagonal of \(e^{z\Delta_{H^4}+w\Delta_{H^3}}\) is the \(H^3\) integral of the product of the two heat kernels. Its expansion gives \(D_0=(4\pi)^{-2}z^{-1/2}(z+w\cosh^2\tau)^{-3/2}\), the Gaussian symbol formula above, and, for the pure \(\Box\) term (window removed), the Seeley--DeWitt coefficient \((5m^2-3\ell^2)/(4m^2)\), \(m=y^4\), the formal density of the preceding paragraphs obtained by an independent route. With the window the weight-two part of \(\sum\ln(\lambda-i\epsilon)\) is \(-(4\pi)^{-2}\LUV^2\,Q(a)\), \(a=\sinh\tau\), plus \(-2\) times a scale-dependent logarithm proportional to the number of modes; \(\operatorname{Re}Q\) changes sign at \(a\simeq0.62\) (not at \(a^2=\tfrac35\), the zero of the coefficient above) and is positive above, and \(\operatorname{Im}Q\to\pi\). This is the Lorentzian density of one scalar sector at the background \(A=0\), with a scale-dependent additive constant; it is not the Einstein--Hilbert coefficient.

\emph{Consequences.} (i) The coefficient carrying \(w\) is the weight-two term, a local Poisson-algebraic coefficient multiplied by window moments; the signature enters through the moments and the window, and in the potential channel the Lorentzian moment is finite and positive in the class sampled. (ii) The mode-count term is proportional to the difference between the numbers of bosonic and fermionic modes in the window, so its logarithmic scale dependence cancels only if that difference vanishes at this order. (iii) The weight-two coefficient as a quadratic form in \(A\) is not computed, the vector, ghost and fibre-harmonic sectors are not computed, and neither the sign nor the size of \(w\), hence no sign of an attractive coupling, follows.

This is a statement for the Poisson limit, one loop, the leading symbol and the invariant sector, the four-matrix class with \(\mu^2=-3\), Gaussian windows and the sampled points; the admissible cone is conditional on the identification stated, and the positivity of the moment is numerical. It does not treat fermions, Moyal terms, the dependence on \(A\), or a window built from dynamical matrices only; it identifies no induced metric with the emergent metric of Chapter~\ref{ch:vierbein-emergence} and, since a Euclidean coefficient is not a Lorentzian one, leaves Assumption~\ref{ass:lorentzian} in place (Problem~E). This changes no ledger label; Problem~F remains open.
